% !PRINTED 315
\weberpage{315}
%<<CONT
divisors we have already considered (\S.\,70 to 72), and all proper divisors of $S$ are metacyclic. Among them are divisors of degree $21$ and $7$, which we shall denote by $T_{21}$, $T_7$, further divisors of degree $24$, and still other divisors whose degree numbers divide 24, all of which we shall here comprise under the symbol $T_{24}$.

These groups $T_{24}$ cannot connect the seven quantities $(\xi) = (\xi_1, \xi_2, \xi_3)$ transitively with one another, because the degree of a transitive permutation group of seven elements must be divisible by 7 (Bd.\,I, \S.\,154, 7.). We shall show that the permutation groups of the $X$ corresponding to the groups $T_{24}\,Q$ are imprimitive. % [sic] Full stop before the syntactic continuation with ``and that.''
and that these groups therefore drop out of our consideration\footnote{From the literature on the equations of degree $8$ the following are to be mentioned: \gesperrt{Nöther}, Mathem. Annalen, Bd. XV. \gesperrt{Wilshaus}, Ueber die algebraische Auflösbarkeit der Gleichungen $8^{\mathrm{ten}}$ Grades. Dissertation. Marburg 1888.}.

Since all permutations $\sigma$ which leave two of the seven elements $(\xi)$ unchanged, or merely permute them among themselves, leave a third element unchanged, namely the complement to the triple, two cases are possible:

\begin{itemize}
\item[1)] the group $T_{24}$ leaves one element at rest, or
\item[2)] the seven elements are decomposed into two parts of three and four elements, of which each part is permuted only among itself.
\end{itemize}

In the latter case we may regard the three elements as belonging to a triple, and accordingly the four elements as belonging to a quadruple. For if three elements which do not belong to a triple are permuted among themselves, then the three complements of each pair of elements to the triple are also permuted among themselves, and the one remaining element stays unchanged. We thus come back to case 1).

We now add to the seven systems $(\xi)$ the eighth $(0,0,0) = (0)$ and denote the eight quantities $X_{\xi_1, \xi_2, \xi_3}$ by

\weq{1}{0, 1, 2, 3, 4, 5, 6, 7.}

In case 1) two of these eight elements, say 0, 1, then remain unchanged under the substitutions of $T_{24}$. There is, further, in $Q$ one and only one substitution $\gamma_0$ by which 0 is interchanged with 1.
